\section{Resampling a spread family}\label{sec:resampling}

The graph argument will produce distributions of edge extensions. We first
show how one product sample can replace a set from such a distribution
by a smaller remainder, and bound the resulting change in the law of
coordinates outside the chosen set.

For $a>0$, a probability distribution on subsets $A\subseteq X$ is
\emph{$a$-spread} if
\begin{equation}\label{eq:spread}
 \PP(J\subseteq A)\le a^{|J|}
 \qquad(\varnothing\ne J\subseteq X).
\end{equation}
We allow the distribution to be carried by a finite indexed family
$(A_b)_{b\in B}$: distinct indices may have the same label, and labels
may be empty. Write $\nu_b>0$ for the weights, with $\sum_b\nu_b=1$.

The transition below is the planted-set and posterior-resampling
construction of Mossel, Niles-Weed, Sun, and Zadik~\cite[Sections~2.1--2.3]{MNSZ25}.
We record a capacity-dependent failure estimate and a change-of-measure
bound for functions of coordinates outside the target label; the latter
will control the dependence between successive tree levels.

Fix $0<\rho<1$ and $W\subseteq X$. The following transition will be
used at every node of a tree:
\begin{enumerate}[label=\arabic*.,leftmargin=2em,itemsep=2pt]
 \item Choose a \emph{source} $c$ with law $\nu$, and set $Y=W\cup A_c$.
 \item Choose a \emph{target} $b$ with conditional probabilities
 \begin{equation}\label{eq:posterior}
  \frac{\nu_b\rho^{-|A_b|}\ind_{\{A_b\subseteq Y\}}}{Z(Y)},
  \qquad
  Z(Y)=\sum_{d\in B}\nu_d\rho^{-|A_d|}
                  \ind_{\{A_d\subseteq Y\}}.
 \end{equation}
 \item Retain the fragment $T=A_b\setminus W$.
\end{enumerate}
The denominator is positive because $A_c\subseteq Y$ and $\nu_c>0$.
To see why these are posterior probabilities, now let $W\sim\mu_\rho$
independently of the source. For a fixed source $c$, the probability
mass of the planted union at a set $Y$ is
$\mu_\rho(Y)\rho^{-|A_c|}\ind_{\{A_c\subseteq Y\}}$.
After averaging over $c$, the law of $Y$ therefore has density $Z$
relative to $\mu_\rho$, and Bayes' formula gives
\eqref{eq:posterior}. At a specified $W$, however, only the source
is guaranteed to have its original law $\nu$; the target law may differ.

For every transition of positive probability,
\begin{equation}\label{eq:two-inclusions}
 T\subseteq A_c\cap A_b,
 \qquad A_b\subseteq W\cup T.
\end{equation}
The first inclusion will preserve spread by comparing the fragment
with the source. The second lets us recover the chosen target from
$W$ and its fragment. In the tree application, that target also
specifies which child and its continuations are kept.
The entire transition, including its fragment, only inspects
membership in $W$ on $\bigcup_{b\in B}A_b$.

\begin{lemma}[Local resampling]\label{lem:local}
Suppose that $|A_b|\le m$ for a positive integer $m$, the law $\nu$
is $a$-spread, and
\[
 0<\rho<1,\qquad a/\rho\le\ee^{-50}.
\]
Let $W\sim\mu_\rho$, independently of the source choice, and perform
the transition above. Call the transition a \emph{local failure} if
\[
 Z(Y)<\ee^{-10m}\quad\text{or}\quad |T|>m/2.
\]
Then
\begin{equation}\label{eq:local-failure}
 \PP(\text{local failure})\le\ee^{-9m}.
\end{equation}
Moreover, for each fixed target $b$ and every nonnegative function
$f$ depending only on $W\cap(X\setminus A_b)$,
\begin{equation}\label{eq:distortion}
 \EE\!\left[
  \ind_{\{\mathrm{target}=b,\ Z(Y)\ge\ee^{-10m}\}} f(W)
 \right]
 \le \ee^{11m}\nu_b\,\EE f(W).
\end{equation}
The expectation on the right is with respect to $W\sim\mu_\rho$.
\end{lemma}

\begin{proof}
The density identity above gives
\begin{equation}\label{eq:small-denominator}
 \PP\bigl(Z(Y)<\ee^{-10m}\bigr)
 =\EE_{Y\sim\mu_\rho}
   \bigl[Z(Y)\ind_{\{Z(Y)<\ee^{-10m}\}}\bigr]
 \le\ee^{-10m}.
\end{equation}

We will use the following consequence of spread. For a fixed $b$,
let $R=|A_c\cap A_b|$, where $c$ has law $\nu$. For every
$0\le r\le m$, including the trivial case $r=0$,
\begin{equation}\label{eq:intersection}
 \PP(R=r)
 \le\sum_{\substack{J\subseteq A_b\\|J|=r}}
        \PP(J\subseteq A_c)
 \le\binom mr a^r.
\end{equation}
On $Z(Y)\ge\ee^{-10m}$, bound the reciprocal denominator in
\eqref{eq:posterior} by $\ee^{10m}$. For a fixed pair $(c,b)$,
the condition $A_b\subseteq W\cup A_c$ requires
$A_b\setminus A_c\subseteq W$, an event of probability
$\rho^{|A_b\setminus A_c|}$. Also, $|T|>m/2$ implies
$|A_c\cap A_b|>m/2$. Therefore
\begin{align*}
 &\PP\bigl(Z(Y)\ge\ee^{-10m},\ |T|>m/2\bigr)\\
 &\qquad\le
 \ee^{10m}\sum_{c,b}\nu_c\nu_b
 \rho^{-|A_c\cap A_b|}
 \ind_{\{|A_c\cap A_b|>m/2\}}\\
 &\qquad\le
 \ee^{10m}\sum_{r>m/2}\binom mr(a/\rho)^r
 \le \ee^{10m}2^m\ee^{-25m}.
\end{align*}
Together with \eqref{eq:small-denominator}, this gives
\[
 \PP(\text{local failure})
 \le\ee^{-10m}+\ee^{-(15-\log 2)m}
 \le\ee^{-9m},
\]
as $m\ge1$.

For \eqref{eq:distortion}, again bound the reciprocal denominator
and then sum over sources. The event
$A_b\setminus A_c\subseteq W$ uses only coordinates inside $A_b$,
so it is independent of $f(W)$. Hence
\begin{align*}
 &\EE\!\left[
 \ind_{\{\mathrm{target}=b,\ Z(Y)\ge\ee^{-10m}\}} f(W)
 \right]\\
 &\qquad\le
 \ee^{10m}\nu_b\sum_c\nu_c\rho^{-|A_b|}
       \EE\bigl[\ind_{\{A_b\setminus A_c\subseteq W\}}f(W)\bigr]\\
 &\qquad=
 \ee^{10m}\nu_b\,\EE f(W)
       \sum_c\nu_c\rho^{-|A_c\cap A_b|}\\
 &\qquad\le
 \ee^{10m}\nu_b\,\EE f(W)
       \sum_{r=0}^m\binom mr(a/\rho)^r
 \le\ee^{11m}\nu_b\,\EE f(W).
\end{align*}
The penultimate inequality is \eqref{eq:intersection}; the last uses
$(1+a/\rho)^m\le\ee^m$.
\end{proof}
